%% kapa: exam=1 round=22 year=2011 period=2 form=A subject=accounting subject_ko=회계학 seq=10 number=10
%% source: Q-Net 감정평가사 기출문제 (cst00309) artlSeq=5036790 file=extracted/1008/5036790/2011년 제22회 감정평가사 1차시험 기출문제/2교시.hwp
%% conversion: kordoc 4.24.51 -> tools/gen1.py (auto); review=proofread-gpt-astra; flags=none
\begin{question}{10}
기능통화에 의한 외화거래의 인식 및 측정으로 옳지 않은 것은?
\begin{choices}
  \item 기능통화로 외화거래를 최초로 인식하는 경우에 거래일의 외화와 기능통화 사이의 현물환율을 외화금액에 적용하여 기록한다.
  \item 거래일은 거래의 인식조건을 최초로 충족하는 날이다. 실무적으로는 거래일의 실제 환율에 근접한 환율을 자주 사용한다.
  \item 공정가치로 측정하는 비화폐성 외화항목은 평균환율로 환산한다.
  \item 역사적원가로 측정하는 비화폐성 외화항목은 거래일의 환율로 환산한다.
  \item 비화폐성항목에서 생긴 손익을 기타포괄손익으로 인식하는 경우에 그 손익에 포함된 환율변동효과도 기타포괄손익으로 인식한다.
\end{choices}
\end{question}
\answer{3}
